\subsection{Ergebnisse der Umsetzung}
Durch die Umsetzung gibt es ein k3s-Cluster mit HA-Architektur für die Produktionsumgebung im Hochschulnetzwerk. 
Da die Anwendung im Hochschulnetzwerk bleibt, müssen die VMs nicht über Tools wie OpenTofu und Proxmox selbst verwaltet werden, sondern weiterhin von der IT.
Der Cluster kann per SSH auf einen Master-Node zugegriffen werden. Es gibt aber auch die Möglichkeit, einen Admin-Rechner einzurichten. Dieser kann über den LB auf den Cluster zugreifen. 
Das ist praktisch, wenn es benötigt wird, bei einen größeren Betrieb.

Diese Ergebnisse steigern die Zuverlässigkeit. 
Dank der HA-Architektur hat die Anwendung nicht nur eine höhere Verfügbarkeit, sondern auch eine bessere Fehlertoleranz. So können nicht nur Ausfälle von Worker Nodes, sondern auch von Master Nodes toleriert werden.
Durch den Load Balancer und die mehreren Master- und Worker-Nodes steigt auch die Stabilität der Anwendung, da diese ihre Last nun besser verteilen kann.

Hinzu kommt, dass dieses System skalierbar ist. Wenn bei einen größeren Betriebsverhalten mehr Last erzeugt wird, können weitere Master- und Worker-Nodes zum Cluster hinzugefügt werden, um diese besser zuverteilen.
Auch der LB kann Replicas bekommen bei einen Wachstum der Nodes. Daher kann der Cluster mit Codetask zusammen wachsen.

